% Аннотация (русская и английская) и заявление об использовании ИИ. Дописать после пилота:
% предложение о результатах помечено \pending, пока у RQ-1…RQ-3 нет чисел с паспортом (docs/thesis/figures.md).
\chapter*{Аннотация}
\addcontentsline{toc}{chapter}{Аннотация}
\thispagestyle{plain}

Взрослые, которые учатся читать Коран, обычно начинают с букваря и учителя, слушающего каждый слог.
В диссертации исследуется, какую часть этого слушания может взять на себя машина для русскоязычных
начинающих, и создаётся инструмент для ответа на этот вопрос. Инструмент --- Telegram-тренажёр, который
ведёт ученика по классическому букварю Ахмад-Хади Максуди «Муаллим ас-Сани» в порядке самой книги:
125 уроков и 800 карточек, каждая прослеживается до страницы и строки издания 2010~г. Ученик присылает
голосовое; открытая модель Quran Muaalem распознаёт фонемы; выравнивание по правилам относит расхождения
к десяти типам ошибок с привязкой к слогам; ученик получает короткое объяснение на русском и эталонную
запись. По каждой попытке сохраняется полный след решения, поэтому любой вердикт можно пересчитать.

Работа выполнена в парадигме design science и оценивает два артефакта. Тренажёр оценивается
измерением согласия с двумя преподавателями (RQ-1), анализом ошибок начинающих, в котором ошибки
ученика отделены от ошибок модели (RQ-2), и внутрисубъектным пилотным экспериментом, сравнивающим полные
объяснения с сигналом «верно / неверно» (RQ-3). Обвязка разработки --- хуки, которые замораживают тесты,
защищают зоны человека и не дают агенту закончить при красных проверках, вместе с
субагентами-ревьюерами и журналами --- оценивается как встроенное кейс-стади (RQ-4). Гипотезы и план
анализа предрегистрированы.

Предварительные проверки показывают, что конвейер отличает верное чтение эксперта от неверного
(99--100 против 40--67 баллов) и отвечает быстрее 10 секунд. \pending{результаты пилота по RQ-1--RQ-3}.
Вклад работы --- первая оценка открытой модели проверки чтения Корана на букварном чтении начинающих,
для которых арабский не родной, относительно оценки преподавателей; набор данных с двойной разметкой,
собранный с согласия участников; доказательное описание разработки исследовательского ПО ИИ-агентами под
механическим контролем.

\medskip
\noindent\textbf{Ключевые слова:} компьютерное обучение произношению; обнаружение ошибок произношения;
чтение Корана; Telegram-бот; design science; ИИ-агенты; обвязка агента.

\chapter*{\textenglish{Abstract}}
\addcontentsline{toc}{chapter}{Abstract (аннотация на английском)}
\thispagestyle{plain}
\begin{otherlanguage}{english}
Adults who learn to read the Qur'an usually start with a primer and a teacher who listens to every
syllable. This thesis asks how much of that listening can be done by a machine for Russian-speaking
beginners, and builds the instrument needed to answer the question: a Telegram tutor that follows the
primer \emph{Muallim al-Sani} by Ahmad Hadi Maksudi in the book's own order, with 125 lessons and 800 cards
traceable to a page and line of the 2010 edition. The learner sends a voice message; the open model Quran
Muaalem recognises phonemes; a rule-based alignment classifies differences into ten error kinds tied to
syllables; and the learner receives a short explanation in Russian with a reference recording.

Following design science research, the thesis evaluates the tutor by agreement with two teachers (RQ-1),
an analysis of beginners' errors that separates learner errors from model errors (RQ-2) and a within-subject
pilot experiment on explanations (RQ-3), and evaluates the AI-agent development harness as an embedded case
study (RQ-4). Preliminary checks show that the pipeline separates correct from incorrect expert recitation
(99--100 versus 40--67 points). \pending{pilot results}. The contribution is the first teacher-referenced
evaluation of an open Qur'anic pronunciation model on primer-level reading by second-language beginners, a
consented, doubly annotated dataset, and an evidence-based account of building research software with AI
coding agents under mechanical guard-rails.

\medskip
\noindent\textbf{Keywords:} computer-assisted pronunciation training; mispronunciation detection;
Qur'anic Arabic; Telegram bot; design science; AI coding agents; harness engineering.
\end{otherlanguage}

\chapter*{Заявление об использовании ИИ-инструментов}
\addcontentsline{toc}{chapter}{Заявление об использовании ИИ-инструментов}
\thispagestyle{plain}

Программное обеспечение, описанное в диссертации, разработано автором вместе с ИИ-агентом (Claude Code).
Агент написал большую часть кода и черновиков документации по правилам, описанным в
разделе~\ref{sec:impl-harness}; каждый тест принят автором до реализации, а все решения об учебном
контенте, текстах для ученика, порогах, согласии и данных принимал автор. Черновики этого текста
подготовлены тем же агентом по материалам проекта и отредактированы автором. Каждый запрос к агенту
сохранён с датой в журнале проекта, поэтому использование ИИ можно проверить. Агент не размечал данные
и не вычислял ни одного результата вне версионированного кода, на который ссылается каждая таблица.
Библиографические данные сверены с Crossref, arXiv или сайтом издателя. Заявление будет приведено в
соответствие с политикой университета об использовании ИИ, когда она будет подтверждена.
